\clearpage
\section{Terminology and Compatibility}\label{page:terminology-and-compatibility}\label{section:terminology}

\subsection{Address Values}\label{section:address-terminology}

An \emph{effective address (EA)} is the address value or operand designator produced by an addressing mode before
segment pre-translation. A memory EA produces an address, while register and immediate EAs name non-memory operands.

A \emph{pre-segment virtual address} is a virtual-address value before segment pre-translation. Segment
pre-translation either converts it to a linear address or reports the architecturally defined fault.

A \emph{linear address} is the result of segment pre-translation. The page-table walker consumes it when paging is
enabled; otherwise the memory system uses it directly.

A \emph{physical address} is the byte address produced by page-table translation for ordinary byte-addressed memory.
With paging disabled, the linear address is used directly instead.

A \emph{memory-system address} is the value presented to the memory subsystem after every active architectural
translation stage. Depending on the final access attributes, the subsystem may interpret it as an ordinary byte
address or an externally acknowledged slot address.

\emph{Byte-addressed normal memory} is the coherent memory-location model selected by a final translation with
\texttt{AT=0}. The PTE \texttt{CP} field selects one of its four normal-memory cache policies. Only policy
\texttt{CP=0} is \emph{cacheable}; the other three \texttt{CP} values remain byte-addressed normal memory with the
allocation and propagation behavior defined by the Memory Model.

A \emph{slot-addressed transaction} is the externally acknowledged access selected by a final translation with
\texttt{AT=1}. Its memory-system address is a slot address rather than a normal-memory byte address. The
Memory Address Translation and Memory Model sections define its transfer, completion, and ordering rules.

\emph{Address generation} is the instruction-local calculation that combines registers, displacement, index,
scale, and immediate fields into an effective address. It does not by itself read memory.

\subsection{Segmentation and Paging}

\emph{Segment pre-translation} is the stage between effective-address calculation and paging. A disabled segment
passes the EA through. A translated window checks an offset and adds its base, while a bounds-only window checks an
already-linear address without adding the base.

A \emph{segment register image} is the 64-bit architectural value containing the segment base page, size exponent,
size mantissa, and bounds-only enable bit.

A \emph{disabled segment} has a zero size mantissa. It performs no segment-window check or base addition and passes
the effective address through as the linear address.

A \emph{translated segment} is enabled with bounds-only mode clear. It checks the effective address as an offset
within the segment span and then adds the segment base.

A \emph{bounds-only segment} is enabled with bounds-only mode set. It treats the effective address as already linear
and checks only that it lies inside the segment window.

\emph{Page-table translation} is the optional \texttt{PTCR.PE}-controlled stage that maps a canonical linear
address to a physical frame and applies PTE permissions, cache policy, and addressing type.

A \emph{PFN} is a physical frame number. A byte-addressed leaf PTE at level $L$ combines its naturally aligned PFN
base with the low $12+9(L-1)$ linear-address bits. An L1 slot leaf combines its PFN with linear-address bits 11..0.

\subsection{Control Flow}\label{section:control-flow-terminology}

A \emph{near control transfer} is a CALL, RET, JMP, or conditional branch that changes only PC within the current
code-segment context.

A \emph{segment-changing control transfer} updates CS and PC in one architectural commit. Long jumps, long calls,
and long returns belong to this category.

\subsection{Tracing Terms}\label{section:tracing-terminology}

The \emph{\texttt{TRACE} instruction} records its immediate marker and current instruction boundary in the
implementation trace stream when that stream is enabled. A \emph{trace unit} is the execution-completion unit
sampled by \texttt{STATUS.TF}. \emph{\texttt{DEBUG\_TRACE}} is the architectural event delivered after successful
completion of a TF-selected trace unit. These names identify the marker instruction, execution boundary, and event,
respectively.

\subsection{Architectural State}\label{section:architectural-state-terminology}

A \emph{general-purpose register} is one of the sixteen 64-bit Rn integer registers, R0 through R15. Ordinary
integer register fields use four bits to select R0 through R15.

An \emph{Rn register} is a general-purpose integer register selected by a four-bit Rn field. It may hold integers,
addresses, counters, selectors, or temporary data according to the instruction form.

The \emph{stack pointer register} is SP. It lies outside the four-bit Rn namespace and uses SS as the fixed segment
for SP-relative addressing.

A \emph{segment register operand} is a 64-bit segment register exposed through the S operand class. The SREG selector
table maps its eight encodings to DS, SS, and GS0 through GS5. \texttt{RDSEG CS} and \texttt{PUSH CS} read the
current CS image through their fixed CS forms.

A \emph{state register} is named architectural state such as FLAGS or STATUS. \texttt{RDFLAGS}, \texttt{WRFLAGS},
\texttt{RDSTATUS}, and \texttt{WRSTATUS} access these registers.

A \emph{control register} is a privileged register exposed through the CR operand class, such as PTCR, ASCR, ECR,
SPC, SCS, and SDS.

\subsection{Compatibility}\label{section:compatibility-terminology}

\emph{Reserved} means that software must not depend on the field's value or behavior. \emph{Must be zero} requires
software to write zero, and \emph{read-as-zero} means reads return zero while the field remains reserved.
\emph{Ignored} means hardware accepts the field without using it for the current operation. \emph{Preserved}
requires software to retain the previous value when rewriting the containing object.

\emph{Software-defined} state is reserved for software use; hardware interpretation is limited to behavior
explicitly defined by the containing structure. \emph{Implementation-defined} behavior is selected by the
implementation and published through CPUID, platform documentation, or firmware. \emph{Illegal} use of an opcode,
effective-address form, selector, or operation raises the exception named by the defining rule.
